\section{Convergence, noise and the limits of learning}\label{ch:convergence}
\sectionbrief
{one analytic function $u$ on a connected open set $\Omega\subset\R^n$, observed through its Taylor jet at a point $p_0$ (the datum at step $K$ is $j_Ku$), or through a noisy copy $\tilde u$ with $\|\tilde u-u\|\le\delta$ in a stated norm.}
{Sections~\ref{ch:ident} and~\ref{ch:data} proved that no finite amount of data certifies identification. A real experiment gives more and more data: does the identified operator converge, how fast, and what can a noisy measurement still say?}
{supplies Steps~1--4 and~7 of the extended algorithm (Section~\ref{ch:extalgo}). The proofs are Dim.~A41--A43.}

\subsection{Identification in the limit}

At step $K$ the rule computes, for each degree $m\le K$, the space $N_{m,K}$ of operators of degree $\le m$ whose action on $u$ has vanishing Taylor coefficients up to order $K-m$. These spaces decrease with $K$ and their limit is $\Ann(u)\cap R_{\le m}$. The \emph{minimal-degree rule} outputs the gcd of $N_{m^*(K),K}$ at the smallest degree $m^*(K)$ where it is non-zero.

\begin{theorem}[{Identification in the limit \led{CONV-L1, CONV-T1} \proved}]\label{thm:convt1}
If $\Ann(u)\ne0$ with minimal degree $d$, the rule outputs $g_d=\gcd(\Ann(u)\cap R_{\le d})$ from some step on; if $\Ann(u)=(P)$ it outputs $P$, without knowing $\deg P$; if $\Ann(u)=0$ the degree $m^*(K)\to\infty$. The sequence of outputs takes at most $(d+1)(d+2)/2$ values. But:
\begin{itemize}
\item the step $K^*$ after which the output is correct is not bounded by any function of $d$;
\item no function of the jet of order $K$ is a correct certificate: any finite jet is shared by a polynomial, whose annihilator is a different ideal.
\end{itemize}
\dimref{CONV-T1}
\end{theorem}

\begin{example}{Convergence at an unpredictable time}
For $u=1+x^{N+2}e^{x^2}$ one has $\Ann(u)=(\tau)$. The rule outputs ``no law of degree $1$'' for $K\le N+1$ and $\tau$ from $K=N+2$, for every $N$: nothing in the first $N+1$ derivatives announces the law.
\end{example}

\subsection{No learner decides whether a law exists}

\begin{theorem}[{Analytic Gold theorem \led{CONV-T2, CONV-T2$'$} \proved}]\label{thm:convt2}
For any sequence of maps $\varphi_K$ from jets of order $\le K$ to $\{\text{law},\text{no law}\}$ (no computability or measurability assumed) there is a real entire function $u$ on which $\varphi_K(j_Ku)$ does not converge to the correct answer. The same holds restricted to the class where $\Ann(u)$ is principal or zero; hence no learner identifies $\Ann(u)$ on that class.
\dimref{CONV-T2}
\end{theorem}

\begin{intuition}
Build $u$ as a limit of polynomials. Whenever the learner says ``law'' on the current polynomial (it must, eventually), add Taylor coefficients, chosen outside a proper algebraic set, that destroy every law of degree $\le k$. The limit function is differentially transcendental yet agrees with a polynomial where the learner said ``law'', infinitely often.
\end{intuition}

\begin{keypoint}
``$\Ann(u)\ne0$'' is \emph{verifiable in the limit} (a rule that is eventually right when it holds) but not \emph{decidable in the limit}.
\end{keypoint}

\subsection{Noisy data}

\begin{theorem}[{Noise: discontinuity, conditional stability, discrepancy rule \led{CONV-T3, T4, T5} \proved}]\label{thm:convt45}
(T3) In spaces where polynomials are dense and $e^{|z|^2}$ lies, the sets $\{\Ann=0\}$ and $\{\Ann\ne0\}$ are both dense: $u\mapsto\Ann(u)$ is discontinuous everywhere. (T4) At a fixed level $n_0$ with a margin $\rho=\sigma_2(\Phi(u))>0$, the minimal singular vector of $\Phi(\tilde u)$ is within $\sin\angle\le2\delta/\rho$ of the true law, an order that is optimal; under an $L^2$ error and an a priori $H^{n_0+s}$ bound the rate is H\"older, and $\sigma_1(\Phi(\tilde u))>\delta$ certifies a misfit. (T5) The discrepancy rule $m(\delta)=\min\{m:\sigma_1(\Phi_m(\tilde u))\le\delta_m\}$ returns the correct degree $d$ and converges at rate $\delta$ once $\delta<\delta_0=\tfrac12\min_{m<d}s_m$; if $\Ann(u)=0$, $m(\delta)\to\infty$. The threshold $\delta_0$ cannot be bounded below from the data.
\dimref{CONV-T4}
\end{theorem}

\begin{example}{The heat kernel from noisy data}
For $u=(t+1)^{-1/2}e^{-x^2/(4(t+1))}$ (minimal degree $2$, $\Ann=(\tau-\xi^2)$), the rule returned $m(\delta)=2$ for $\delta=10^{-3}$ and $10^{-4}$ in a numerical check.
\end{example}

\subsection{Other results of this phase}

\begin{itemize}
\item \textbf{Locality}: when the law is a non-local operator it is invisible to $\Ann(u)$; an irrational fractional power or a delay gives $\Ann=0$ (a delay is in the local class through an auxiliary transport) \led{CONV-T6, CONV-P6, CONV-P7}.
\item \textbf{Dimension $1+d$}: the ambiguity of the factor decomposition is infinite unless the factors are pairwise co-maximal \led{CONV-D1}; on a convex region $\ker(Q_1Q_2)=\ker Q_1+\ker Q_2$ in $C^\infty$ \led{CONV-D2}.
\item \textbf{Boundary conditions} are invisible to $\Ann(u)$; stability holds for a dictionary of traces \led{CONV-B1--B3}. Splitting with a Dirichlet boundary has global order $1$ (Lie) and $\Theta(h^{5/4})$ (Strang) in $L^2$ \led{CONV-T8}.
\item \textbf{Mori--Zwanzig} for constant-coefficient $2\times2$ systems: exact memory form and eliminated law \led{CONV-MZ1--3}; the criterion ``$E_{\rm int}=0\iff[\mathcal P,L]=0$'' is false \led{CONV-MZ3-F}.
\item \textbf{A convergence contract} for chains of steps: classes (A) exact, (B) convergent with explicit error, (C) no guarantee, (Lim) in the limit; the error of a chain of Lipschitz steps is $\le\sum_i(\prod_{j>i}L_j)\varepsilon_i$ \led{CONV-T7}.
\end{itemize}

\begin{takeaway}
Learning an operator from growing data is possible only in the limit and without a certificate; deciding whether any law exists is impossible even in the limit; and noise can be handled only relative to a finite level of the dictionary and a margin that the data cannot supply.
\end{takeaway}
